% Options for packages loaded elsewhere
\PassOptionsToPackage{unicode}{hyperref}
\PassOptionsToPackage{hyphens}{url}
\PassOptionsToPackage{dvipsnames,svgnames,x11names}{xcolor}
\documentclass[
  10pt,
  fontset=fandol]{ctexart}
\usepackage{xcolor}
\usepackage[margin=16mm]{geometry}
\usepackage{amsmath,amssymb}
\setcounter{secnumdepth}{-\maxdimen} % remove section numbering
\usepackage{iftex}
\ifPDFTeX
  \usepackage[T1]{fontenc}
  \usepackage[utf8]{inputenc}
  \usepackage{textcomp} % provide euro and other symbols
\else % if luatex or xetex
  \usepackage{unicode-math} % this also loads fontspec
  \defaultfontfeatures{Scale=MatchLowercase}
  \defaultfontfeatures[\rmfamily]{Ligatures=TeX,Scale=1}
\fi
\usepackage{lmodern}
\ifPDFTeX\else
  % xetex/luatex font selection
\fi
% Use upquote if available, for straight quotes in verbatim environments
\IfFileExists{upquote.sty}{\usepackage{upquote}}{}
\IfFileExists{microtype.sty}{% use microtype if available
  \usepackage[]{microtype}
  \UseMicrotypeSet[protrusion]{basicmath} % disable protrusion for tt fonts
}{}
\makeatletter
\@ifundefined{KOMAClassName}{% if non-KOMA class
  \IfFileExists{parskip.sty}{%
    \usepackage{parskip}
  }{% else
    \setlength{\parindent}{0pt}
    \setlength{\parskip}{6pt plus 2pt minus 1pt}}
}{% if KOMA class
  \KOMAoptions{parskip=half}}
\makeatother
\usepackage{longtable,booktabs,array}
\newcounter{none} % for unnumbered tables
\usepackage{calc} % for calculating minipage widths
% Correct order of tables after \paragraph or \subparagraph
\usepackage{etoolbox}
\makeatletter
\patchcmd\longtable{\par}{\if@noskipsec\mbox{}\fi\par}{}{}
\makeatother
% Allow footnotes in longtable head/foot
\IfFileExists{footnotehyper.sty}{\usepackage{footnotehyper}}{\usepackage{footnote}}
\makesavenoteenv{longtable}
\setlength{\emergencystretch}{3em} % prevent overfull lines
\providecommand{\tightlist}{%
  \setlength{\itemsep}{0pt}\setlength{\parskip}{0pt}}
\usepackage{amsmath,amssymb,mathtools,needspace}
\xeCJKDeclareCharClass{CJK}{"2032,"0400->"04FF,"2160->"216B,"2460->"2473,"25A0,"25C7,"25CB,"25CF}
\IfFontExistsTF{Arial Unicode MS}{\xeCJKsetup{AutoFallBack=true}\setCJKfallbackfamilyfont{\CJKrmdefault}{Arial Unicode MS}}{}
\setcounter{secnumdepth}{0}
\setlength{\emergencystretch}{2em}
\allowdisplaybreaks
\usepackage{bookmark}
\IfFileExists{xurl.sty}{\usepackage{xurl}}{} % add URL line breaks if available
\urlstyle{same}
\hypersetup{
  pdftitle={Adams 隐式格式},
  colorlinks=true,
  linkcolor={Maroon},
  filecolor={Maroon},
  citecolor={Blue},
  urlcolor={Blue},
  pdfcreator={LaTeX via pandoc}}

\title{Adams 隐式格式}
\author{}
\date{}

\begin{document}
\maketitle

\subsubsection{5.5.3 Adams 隐式格式}\label{adams-ux9690ux5f0fux683cux5f0f}

上述 Adams 显式方法选 \(x_n,x_{n-1},\cdots,x_{n-r}\) 作为插值节点，这时，用插值函数 \(P_r(x)\) 在求积区间 \([x_n,x_{n+1}]\) 上逼近 \(f(x,y(x))\)，实际上是个外推过程，效果不够理想。为了改善逼近效果，可以变外推为内插，改用 \(x_{n+1},x_n,\cdots,x_{n-r+1}\) 为插值节点；而通过数据点 \((x_{n+1},\allowbreak f_{n+1}),\allowbreak (x_n,\allowbreak f_n),\allowbreak \cdots,\allowbreak (x_{n-r+1},\allowbreak f_{n-r+1})\) 插出函数 \(P_r(x)\)，然后重复前一段的推导过程。对应于式（5.5.3），这里有下列 \textbf{Adams 隐式格式}：

\[
y_{n+1}=y_n+h\sum_{j=0}^r\alpha_j^*\Delta^j f_{n-j+1},\qquad
\alpha_j^*=(-1)^j\int_{-1}^0\binom{-t}{j}\,\mathrm dt.
\tag{5.5.7}
\]

它的前几个值如表 5.8 所示。

\textbf{表 5.8}

{\def\LTcaptype{none} % do not increment counter
\begin{longtable}[]{@{}lrrrr@{}}
\toprule\noalign{}
\(j\) & 0 & 1 & 2 & 3 \\
\midrule\noalign{}
\endhead
\bottomrule\noalign{}
\endlastfoot
\(\alpha_j^*\) & 1 & \(-1/2\) & \(-1/12\) & \(-1/24\) \\
\end{longtable}
}

\textbf{表 5.9}

{\def\LTcaptype{none} % do not increment counter
\begin{longtable}[]{@{}lrrrr@{}}
\toprule\noalign{}
\(i\) & 0 & 1 & 2 & 3 \\
\midrule\noalign{}
\endhead
\bottomrule\noalign{}
\endlastfoot
\(\beta_{0i}^*\) & 1 & & & \\
\(\beta_{1i}^*\) & \(1/2\) & \(1/2\) & & \\
\(\beta_{2i}^*\) & \(5/12\) & \(8/12\) & \(-1/12\) & \\
\(\beta_{3i}^*\) & \(9/24\) & \(19/24\) & \(-5/24\) & \(1/24\) \\
\end{longtable}
}

将差分展开，式（5.5.7）可改写为

\[
y_{n+1}=y_n+h\sum_{i=0}^r\beta_{ri}^*f_{n-i+1},\qquad
\beta_{ri}^*=(-1)^i\sum_{j=i}^r\binom ji\alpha_j^*.
\tag{5.5.8}
\]

表 5.9 提供了系数 \(\beta_{ri}^*\) 的部分数值。式（5.5.8）是含有参数 \(r\) 的一族格式，\(r+1\) 为格式的步数。特别地，一步隐式 Adams 格式（\(r=0\)）为后退的 Euler 格式。两步隐式 Adams 格式（\(r=1\)）为梯形格式，而四步隐式 Adams 格式（\(r=3\)）则为

\[
y_{n+1}=y_n+\frac h{24}(9f_{n+1}+19f_n-5f_{n-1}+f_{n-2}).
\tag{5.5.9}
\]

\end{document}
